%!TEX root = ../hutbthesis_main.tex
\chapter{相关基础理论框架}

\section{CarlaUE4仿真平台基础}

本文使用的是HUTB的人车模拟器，它是为面向具身智能、无人车、无人机而设计的一款高保真的开源仿真系统，用CarlaDReyeVR和OpenSim的原生框架来开发，严格遵照了CarlaUE4核心技术规范和API标准，并对原型进行了一些本土化的改造以及一些能力增强的实现工作，在国内自动驾驶及具身智能方面的研究人员当中得到了较为广泛的应用。该平台弥补了Carla原生版在多智能体协同仿真、中文本地化、多物理场耦合等方面所存在的不足，在一个安全可控的虚拟环境中可以复现现实中不能涵盖到的各种极端工况、长尾情况和危险的交通场景，给自动驾驶感知、决策、规划、控制等各方面的模块测试赋予了一个统一标准化的并且高度适配的实际道路条件下的虚拟仿真实验环境。

\subsection{CarlaUE4平台特性与场景搭建}

hutb 搭建的CarlaUE4 仿真环境，在完全继承原生 Carla 平台核心能力的基础上，进行了针对性的功能扩展与性能优化，其核心特性可归纳为以下 6 个维度：

\subsubsection{多智能体一体化仿真能力}

平台冲破了原生Carla以车辆为界的仿真的束缚，它能支持起具身人、无人车、无人机这三种智能体共同在一个环境中开展联合仿真的工作，而且里面也包含着带有骨骼动画的人行生物力学模子以及无人机飞控动力学模子，可以完成人与车互相配合的情形下展开的各种情况模拟实验，在各种复杂的交通参与者环境之下也能满足许多事情上的跟踪算法去经受考验的需求。同时也对多智能体行为调度逻辑进行了优化，在原来的平台上出现大量的NPC并行运行的时候会出现性能衰减的现象，但是经过处理之后，在一个场景里面可以支撑起上千甚至上万的仿真参与度了。

\subsubsection{本土化高保真开源数字资产库}

平台使用UE4引擎创建了具有影视级3D场景渲染和物理模拟能力的系统，能支持Carla所有的数字资产，并且有自己的开放式非国标开放数字资产管理库（CARTA），包含了按照中国道路标准建好的城镇规划图、建筑物、交通标志标识线以及20种具备完整的动力学参数的国产车、187个带有骨骼动作的行人模型及各种各样的道路器材\cite{zhao2025carlascene}。所有的资产都遵守了统一的物理尺寸的规定，使用git-lfs来处理大型文件资产，并且为了提高渲染速度以及视觉效果，在保证准确性的前提下尽量保持大小不变，可以直接应用于国内交通场景的算法测试当中\cite{li2025virtual}。

\subsubsection{全栈式传感器仿真与扩展能力}

平台完全兼容原生Carla 的全品类传感器模型，支持 RGB 相机、深度相机、语义分割相机、多线激光雷达、毫米波雷达、GNSS、IMU 等传感器的参数化配置，可自定义传感器的安装位置、朝向、视场角、噪声模型、采样频率等关键参数，同时支持地面真值数据的同步输出。在此基础上，平台增加了法线信息传感器等专门用作飞机上空自感知设备的传感器，并且在无人机飞行的时候使用更加方便并且可以满足更多要求的功能也包括进去，在这个系统中包含了激光雷达以及其它一些其他的组件都可以被接入到系统中来，进而给用户增添了一些有用的信息来源。

\subsubsection{低门槛开发与中文本地化适配}

平台对Windows 10+、Ubuntu 18.04+、MacOS 12+全平台都进行了支持，并且对于国内开发者做出了优化，在这里提到了预编译的编辑器版本，只需要点击一次按钮就能直接运行带有插件的虚幻编辑器，以及Python工具包“hutb”，它使用的是Python 3.7到3.14的所有版本，并且通过pip命令可以一次性完成安装过程。同时提供了一个全面的中文搭建文档，开发教程和使用说明，并且附带了多个样例代码。

\subsubsection{交通流与场景引擎深度优化}

平台针对Carla原生的核心功能模块进行了大量的bug修复和性能优化：重新编写了OpenDriveActor模块，修正了仿真环境中原始平台对于环岛路线以及路口的行驶方式处理出现的问题；改善了TrafficManager模块对于大车在路口处进行通行所遇到的瓶颈之处，并且显著降低了在仿真环境当中长时间运转时产生安全事故的风险率；解决了多个传感器同时被创建或者销毁的时候存在的一些网络连接方面的异常情况和内存泄漏现象，进而提升了整个系统的稳定水平。而且平台本身就可以兼容 ScenarRunner 的场景运行引擎，支持 OpenSCENARIO 标准的场景描述方式，并且可以与 SUMO 以及 PTV Vissim 这类交通仿真工具一起进行联合仿真的工作，对于建立一些比较复杂的交通流量场景有着不错的功效表现。

\subsection{基于hutb的仿真场景搭建流程}

本研究采用的 hutb-CarlaUE4 完整环境，分为预编译版本快速部署与源码编译自定义场景开发两种搭建模式，可适配不同的算法测试需求，核心流程如下：

\subsubsection{预编译版本快速部署流程}

根据一般的算法测试要求，可以直接使用平台上预编译的完整的环境，不需要源码编译，主要过程如下：首先从官方发布的页面下载对应的平台预编译安装包并解压运行启动脚本CarlaUE4.exe（Windows）或者Linux中的CarlaUE4.sh来开启hutb-carlaue4仿真服务器，在终端环境中输入命令"pip install hutb"to一键安装Python API工具包以达到客户端环境的搭建目的，再利用客户端API与本地仿真服务器进行连接(默认IP:127.0.0.1，端口:2000)，然后设定各种各样的环境参数如天气、光照、时间等，并且加载相应的中国城镇道路网地图到仿真中去，在此基础上我们还可以创建出一个目标点的主车及多个传感器套件出来，并对其所有的硬件设备以及采集的数据进行设置和追踪功能的操作加以配置之后，就可以按照所设定好的条件和顺序来进行一系列复杂而又繁重的路况模拟试验任务了，而这一切都是依靠着ScenarioRunner这个应用程序程序实现起来的。

\subsubsection{源码编译与自定义场景全流程开发}

针对不同的测试场景需求，平台提供完整的源代码编译和自定义场景开发流程：第一步用 "git clone" 拉取 hutb 分支下的最新源码，依照官方中文文档对 UE4 引擎进行适配、环境搭建和全编译，最终得到可供自定义开发的仿真服务器；第二步使用 RoadRunner 工具根据 OpenDRIVE 标准绘制符合中国道路规范的道路中心线、车道线、交通标志、信号灯等要素图，保存为 OpenDRIVE 地形文件，并将其提交到本地路径下以供编辑；第三部将地形文件导入 UE4 引擎并配置完成本土化的 3D 场景资源，从指定版本包下载静态物体及环境设置（例如建筑模型），然后设定材质样式以及光源参数组合效果，并给各个碰撞体添加必要的属性；第四部利用 Hutb Python/C++ API 编写场景运行程序，实现不同类型车辆之间的交互作用，譬如如何让汽车跟随行人行动，或怎样把风力影响计算进去，又或是经由雷达获取信息后做出相应的驾驶判断，所有这些过程均要依靠自定义场景来执行，且在过程中涉及到了大量技术问题解决工作内容，比如系统性能提升策略分析以及硬件兼容性优化方法探究等方面研究成果颇有价值，本课题致力于探索各种新颖有趣应用领域方向前景表现状况发展趋向特点特征所在之处，其背后所蕴含价值意义远超传统意义上学术层面要求范畴限制范围内所能触及到之处之外。

\section{Chrono::Vehicle车辆动力学模块化建模体系}

Chrono::Vehicle 模块是用模块化、模板化的思想来实现车辆动力学模型的快速建立和灵活改变，也是 Hutb-CarlaUE4 平台上仿真合成功能的基础部分，在本文中所有的车辆动力学模型都是在此基础上搭建起来的\cite{openhutb2026chrono}。

\subsection{模板化建模核心设计}

Chrono::Vehicle 是一个由各种子系统的实例组成的一个类，用基于模板的方式实现建模的灵活性。模板属于车辆子系统参数化的模型，它用来表示汽车各子系统所具有的特性以及具体的设计结果，其主要的功能有：对各个车型的车身进行设定，对车头和座椅等部位加以限定，同时又会对有关于关节以及力学要素方面的事做出界定，并且还会涵盖有关于设计参数的具体要求，除此之外，在实际应用的过程中还能够体现出各种不同型号汽车上装配相同类型的子系统所产生的共通之处。

以悬架系统为例，双叉臂悬架模板已经预先确定好了上控制臂、下控制臂、转向节、减震器、弹簧这几个主要零部件所处的位置以及相互间的约束关系，用户只需要对硬件坐标系（即位置）、关节的方向、惯性的性质、接触材料属性这些模板参数做出调整，并根据实际情况来设定好其他的参数值之后，就可以利用这些模板创建出和某一特定汽车匹配起来的标准前后悬挂结构图象模子，不需要再启动其它方面的编程工作。

完整的一个车辆移动仿真包含三个辅助系统模板：一个是驾驶员系统，用来给车辆的转向、油门、制动这些控制输入赋予相应的命令，并且可以直接同Hutb-carlaue4平台上所具有的车辆控车指令相联系\cite{fu2026carlaoverview}；第二个便是动力总成系统，它是以发动机以及变速箱的形式存在，跟车辆传动系统相连，在进行驱动力计算的时候能做出精确的判断\cite{shi20254ws}；第三个就是地形系统，它可以创建出刚性的或者是柔软的甚至是颗粒状的地形模样来供轮胎和地面之间展开交互使用。流程图如\ref{2.png}所示，功能如表\ref{Description of the functions in the process}：

\begin{figure}[H] % 可以是h（here），t（top），b（bottom），p（page of floats）
	\centering
	\includegraphics[width=1\textwidth]{2} % 假设图片文件名为car.pdf或car.png等，位于当前工作目录
	\caption{数据流闭环解析} % 图片标题
	\label{2.png} % 用于引用的标签
\end{figure}

\begin{table}[H]
\centering
\caption{流程内功能说明}
\label{Description of the functions in the process}
\begin{tabularx}{\linewidth}{@{} >{\raggedright\arraybackslash}p{2.8cm} >{\raggedright\arraybackslash}X >{\raggedright\arraybackslash}X @{}}
\toprule
图内流程 & 类型 & 功能描述 \\
\midrule
DIRVER & 驾驶员 & 提供所有人工操控输入的主体 \\
POWERTRAIN & 动力传动系统（动力总成） & 发动机、变速箱、传动轴等组成的动力传递系统 \\
VEHICLE & 车辆本体（底盘动力学模型） & 包含车身、悬架、转向、制动等机械结构的核心模型 \\
TIRES & 轮胎模型 & 计算轮胎与地面之间作用力的关键模块 \\
TERRAIN & 地形（路面环境） & 提供路面几何、摩擦、高度等信息的环境模块 \\
Throttle input & 油门输入 & 驾驶员控制动力输出的指令 \\
Steering input & 转向输入 & 驾驶员控制车辆转向的指令 \\
Braking input & 制动输入 & 驾驶员控制车辆减速的指令 \\
Driveshaft torque & 传动轴扭矩 & 动力总成传递给车轮的驱动扭矩 \\
Driveshaft speed & 传动轴转速 & 车轮反馈给动力总成的转速信号 \\
Wheel states & 车轮状态 & 车辆本体传递给轮胎模型的车轮运动状态（位置、速度、滑移率等） \\
Height / Normal & 路面高度/法向 & 地形提供给轮胎模型的路面几何与法向信息 \\
Forces and moments on wheel bodies & 车轮受力与力矩 & 轮胎模型计算出的轮胎对车轮的作用力与力矩，反馈给车辆本体 \\
\bottomrule
\end{tabularx}
\end{table}
\subsection{标准化子系统模板库}

\subsubsection{悬架系统模板}
覆盖了双叉臂、基于距离约束的简化后的双叉臂、多连杆、实心轴、麦弗逊式支柱、半拖曳臂这几种常见的悬架结构\cite{xiong2023semi}，适配乘用车辆、商用车辆等不同车型的建模需求\cite{gu2025rbf}。\ref{3.png}所示：
\begin{figure}[H] % 可以是h（here），t（top），b（bottom），p（page of floats）
	\centering
	\includegraphics[width=1\textwidth]{3} % 假设图片文件名为car.pdf或car.png等，位于当前工作目录
	\caption{悬架系统模板} % 图片标题
	\label{3.png} % 用于引用的标签
\end{figure}


\subsubsection{转向系统模板}

给出转向臂和齿条齿轮这两个主要的转向机构模型，并且与前悬架系统相匹配，在展示车辆的转向动力学特性方面也是比较好的。\ref{4.png}所示：

\begin{figure}[H] % 可以是h（here），t（top），b（bottom），p（page of floats）
	\centering
	\includegraphics[width=1\textwidth]{4} % 假设图片文件名为car.pdf或car.png等，位于当前工作目录
	\caption{转向系统模板} % 图片标题
	\label{4.png} % 用于引用的标签
\end{figure}

\subsubsection{传动系统模板}

支撑着2WD两轮驱动和4WD四轮驱动这两类传动系统的建立，用Chrono特有的建模元件来对传动系统进行细致的仿真，还可以得到一些简单的运动学传动系统模板，在满足不同仿真精度的要求时也可以使用到它。\ref{5.png}所示：

\begin{figure}[H] % 可以是h（here），t（top），b（bottom），p（page of floats）
	\centering
	\includegraphics[width=1\textwidth]{5} % 假设图片文件名为car.pdf或car.png等，位于当前工作目录
	\caption{基于轴的 4WD 传动系统与基于轴的动力传动系统模型的连接。} % 图片标题
	\label{5.png} % 用于引用的标签
\end{figure}

Central Diff.(中央差速器)接收变速箱传来的动力，并分给前后轴，可以使得前后轴有不同的转速来提高行驶的稳定程度。图中的箭头为动力传递的路线。

Central diff.（中央差速器）：接收变速箱输出的动力，分配给前后轴，允许前后轴以不同转速旋转，提升行驶稳定性。
动力传递路径如图\ref{Power transmission path}
\begin{figure}[H] % 可以是h（here），t（top），b（bottom），p（page of floats）
	\centering
	\includegraphics[width=1\textwidth]{Power transmission path} % 假设图片文件名为car.pdf或car.png等，位于当前工作目录
	\caption{动力传递路径图} % 图片标题
	\label{Power transmission path.png} % 用于引用的标签
\end{figure}

动力传动系统模板的核心功能如表\ref{Power transmission system table}所示：
\begin{table}[H]
	\centering
	\caption{动力传动系统表}
	\label{Power transmission system table}
	\begin{tabularx}{\linewidth}{lXX}
		\toprule
		概念                 & 范围     & 核心功能   \\
		\midrule
		Powertrain（动力总成）      & 发动机、液力变矩器、变速箱   & 产生动力 + 调节动力（变速变扭）    \\
		Driveline（传动系统） & 变速箱输出轴之后的所有部件（传动轴、差速器、半轴、轮毂）  & 分配动力 + 传递到车轮    \\
		\bottomrule
	\end{tabularx}
\end{table}

\subsubsection{车轮与制动系统模板}

车轮模板是悬架主轴上加装的那些质量、转动惯量以及其它附加件，它用来对车轮惯性的参数做出准确的规定；制动器模板则运用了基于驾驶员制动输入变化来保持固定扭矩的模型，可以重现车辆制动时的动力学反应\cite{guo2025pid}。

\subsubsection{轮胎模型模板}

给出多个精度等级的轮胎模型，即刚性轮胎模型、Pacejka 和 Fiala 的半经验型轮胎模型和采用绝对节点坐标公式或者是共旋转公式的有限元全可变形轮胎模型，并根据仿真需要来挑选合适的轮胎模型，以达到较好的仿真精确度和计算速度的目的。对于履带式的车辆来说，模块还提供有单销和双销的履带板以及链条、扭矩弹簧和液压悬挂装置（扭振减震），并配有挠度调节的挡轮和油压调正装置，在这些参数均可以得到满足的基础上才能对履带式工程车进行动力学计算。

\section{基于JSON文件的车辆模型实例化方法}

为了降低建模门槛，Chrono::Vehicle提供了一种使用JSON格式输入文件来创建车辆模型的方式，并且不需要写C++代码就可以得到一个完整的车辆系统和子系统的模板实例，这也是本文中关于如何改变车辆模型参数主要的方法之一，在Hutb-CarlaUE4平台上可以很方便地从JSON配置文件中加载并使用。根据JSON文件进行车辆建模时，按照子系统的层次结构组织起来，并且使用顶层JSON文件和子系统引用的JSON文件就可以完成一辆轮式或履带式的车辆。双轴双叉臂悬架轮式车辆顶级JSON规范文件的核心结构如上图所示，该结构对车辆的底盘、车轴、悬架、转向、传动系统等各个部分进行完整的定义，并且可以通过改变json文件中对应的参数来实现对于车辆动力学特性的一种更改：

\subsection{坐标系约定与适配规范}

Chrono::Vehicle模块对各种坐标系都有严格的规定，并且为了车辆子系统模型在各平台上都可以兼容、可以进行仿真和测试，它成为与hutb-carlaue4平台相结合的重要基础之一。

\subsubsection{车辆子系统与整车参考系}
每一个车辆子系统都会根据自身所处的参考系来确定所有的硬点和受力点的位置，在模板里面所有的这些信息都需要给出。整车系统用确定各个子系统相对于车辆参考系的位置、方向来装配、连接各个子系统，车辆参考系默认同底盘子系统的参考系相同。模块严格按照ISO 8855：2011车轴约定，使用右手坐标系，X轴向前进向、Z轴向上、Y轴指向车辆左侧，这个约定和hutb-CarlaUE4平台上车辆坐标系完全一致，并且不需要做任何坐标转换就可以做到两个平台之间车辆状态数据互相传递\cite{zhang2025multiscene}。如\ref{6.png}所示：

\begin{figure}[H] % 可以是h（here），t（top），b（bottom），p（page of floats）
	\centering
	\includegraphics[width=1\textwidth]{6} % 假设图片文件名为car.pdf或car.png等，位于当前工作目录
	\caption{车辆参考系O1以及子系统参考系（O′2和 O′′2表示前后悬架，O3表示转向机构）。} % 图片标题
	\label{6.png} % 用于引用的标签
\end{figure}

\subsubsection{模拟世界坐标系}
Chrono::Vehicle默认的世界坐标系和车辆参考系是一致的，使用的是ISO标准的Z轴向上坐标系，还可以设置为自定义的世界坐标系，在不同的仿真环境下可以有不同的约定。对于使用Y轴向上的仿真是可以方便地进行设置和切换，如模块提供了一种快捷的方式来设定为X轴向前、Y轴向上、Z轴向右的世界坐标系即ChWorldFrame::SetYUP$()$；另外，还具有各种工具函数比如坐标系的转换，高度获取以及平面投影等等都可以在ChWorldFrame类里用到，从而能够满足多种仿真平台之间的兼容性需求\cite{ge2025rl}。
